% Katalog C: die Instanz-Regler (Palette Slot 3).
\ZSFNewColumn
\StartChapter[kat:regler]{C · Regler}

\begin{runintext}
  Die Regler sind das eigentliche Gestaltungsangebot. Alle stehen an
  derselben \ZSFkeyword{defbox} mit demselben Mustertext — was sich
  unterscheidet, ist allein der Regler. Vorbelegungen sind nicht als eigener
  Eintrag geführt, sondern in der Zweckzeile genannt; C-00 ist die Referenz.
\end{runintext}

\begin{defbox}[C-00 · Referenz]
  \Muster
\end{defbox}
\Zweck{Dieselbe Box ohne jede Option. Jeder Eintrag unten weicht in genau
einem Regler davon ab.}

\SubsectionBar[kat:regler-farbe]{Farbe und Gewicht}
\ZSFindex{tone}
\ZSFindex{weight}

\begin{defbox}[C-01 · tone=neutral][tone=neutral]
  \Muster
\end{defbox}
\Zweck{Diese Box gehört nicht zum Kapitelthema — Meta-Hinweis, Konvention,
Legende. Vorbelegung: chapter.}

\begin{defbox}[C-02 · tone=warn][tone=warn]
  \Muster
\end{defbox}
\Zweck{Warn-Farbwelt an einer beliebigen Box, ohne den Namen warnbox.}

\begin{defbox}[C-03 · weight=quiet][weight=quiet]
  \Muster
\end{defbox}
\Zweck{Titel wandert in die erste Inhaltszeile, Kopfbalken weicht einem
linken Akzent. Vorbelegung: loud.}

\begin{defbox}[C-04 · weight + tone][weight=quiet, tone=warn]
  \Muster
\end{defbox}
\Zweck{Zwei Regler zusammen — die leise Fassung nimmt den Ton über ihren
Akzent entgegen, nicht über eine Fläche.}

\SubsectionBar[kat:regler-raum]{Innenabstand}
\ZSFindex{padx}
\ZSFindex{pady}

\begin{defbox}[C-05 · padx=bar][padx=bar]
  \Muster
\end{defbox}
\Zweck{Links mehr Rand, damit der Inhalt den Akzentbalken freihält.
Vorbelegung: normal.}

\begin{defbox}[C-06 · padx=none][padx=none]
  \Muster
\end{defbox}
\Zweck{Die Box polstert waagrecht gar nicht — für Inhalt, der es selbst tut.}

\begin{defbox}[C-07 · pady=tight][pady=tight]
  \Muster
\end{defbox}
\Zweck{Knappe Höhe oben und unten. Vorbelegung: normal.}

\begin{defbox}[C-08 · pady=none][pady=none]
  \Muster
\end{defbox}
\Zweck{Senkrecht ohne jede Polsterung.}

\SubsectionBar[kat:regler-kontur]{Kontur und Ausrichtung}
\ZSFindex{frame}
\ZSFindex{align}

\begin{defbox}[C-09 · frame=strong][frame=strong]
  \Muster
\end{defbox}
\Zweck{Kräftigere Kontur; die Farbe kommt weiterhin aus dem Ton.
Vorbelegung: soft.}

\begin{defbox}[C-10 · frame=hard][frame=hard]
  \Muster
\end{defbox}
\Zweck{Die stärkste Stufe.}

\begin{defbox}[C-11 · frame=none][frame=none]
  \Muster
\end{defbox}
\Zweck{Kein Rahmen, Kopfzeile bleibt — Rahmen und Gewicht sind zwei Fragen.}

\begin{defbox}[C-12 · align=center][align=center]
  \Muster
\end{defbox}
\Zweck{Zentrierter Boxinhalt. Vorbelegung: left, also Flattersatz für
schmale Spalten.}

\SubsectionBar[kat:regler-satz]{Satz und Umbruch}
\ZSFindex{font}
\ZSFindex{bodyparskip}
\ZSFindex{atomic}

\begin{defbox}[C-13 · font=dense][font=dense]
  \Muster
\end{defbox}
\Zweck{Inhalt eine Stufe kleiner, Titel unverändert. Vorbelegung: normal.}

\begin{defbox}[C-14 · bodyparskip=inherit][bodyparskip=inherit]
  \Muster

  Sichtbar wird der Unterschied erst am zweiten Absatz.
\end{defbox}
\Zweck{Absatzabstand aus dem Dokument statt aus der Box. Vorbelegung: box.}

\begin{defbox}[C-15 · atomic][atomic]
  \Muster
\end{defbox}
\Zweck{Bricht nie über eine Spaltengrenze — auch nicht im Pack-Modus. Flag
ohne Wert, im Einzeleintrag nicht sichtbar.}

\SubsectionBar[kat:regler-split]{Zweispaltiger Inhalt}
\ZSFindex{split}
\ZSFindex{splitalign}

\begin{defbox}[C-16 · split=0.4][split=0.4]
  \MusterKurz
  \tcblower
  Die rechte Hälfte. Ohne tcblower landet alles links.
\end{defbox}
\Zweck{Anteil der linken Hälfte; gilt auf jeder Box, nicht nur splitbox.}

\begin{defbox}[C-17 · splitalign=top][split=0.4, splitalign=top]
  \MusterKurz
  \tcblower
  Die rechte Hälfte ist absichtlich länger, damit die Ausrichtung der
  beiden Hälften zueinander sichtbar wird.
\end{defbox}
\Zweck{Hälften an der Oberkante — für Diagramm neben Legende.
Vorbelegung: top.}

\begin{defbox}[C-18 · splitalign=center][split=0.4, splitalign=center]
  \MusterKurz
  \tcblower
  Die rechte Hälfte ist absichtlich länger, damit die Ausrichtung der
  beiden Hälften zueinander sichtbar wird.
\end{defbox}
\Zweck{Hälften mittig — für Formel neben Erklärung.}

\begin{defbox}[C-19 · splitalign=bottom][split=0.4, splitalign=bottom]
  \MusterKurz
  \tcblower
  Die rechte Hälfte ist absichtlich länger, damit die Ausrichtung der
  beiden Hälften zueinander sichtbar wird.
\end{defbox}
\Zweck{Hälften auf Grundlinie.}

\SubsectionBar[kat:regler-eigen]{Regler mit nur einem Besitzer}
\ZSFindex{ordered}
\ZSFindex{grid}

\begin{runintext}
  Vier Regler gehören genau einer Box und brechen auf jeder anderen mit einer
  Meldung. Zwei stehen hier, die Tabellen-Regler in \ZSFref{kat:tabellen},
  \ZSFkeyword*{part} in \ZSFref{kat:module}.
\end{runintext}

\begin{ZSFlist}[C-20 · ordered][ordered]
  \ZSFItem{Erster Schritt}{Mustertext zum Formvergleich.}
  \ZSFItem{Zweiter Schritt}{Die Nummerierung kommt aus dem Regler.}
\end{ZSFlist}
\Zweck{ZSFlist nummeriert statt aufgezählt — nur bei stabiler Schrittfolge.}

\begin{valuegrid}[C-21 · grid=horizontal][grid=horizontal]{2}
  Mustertext & zum & Formvergleich \\
  Zweite Zeile & nur & Zeilenlinien
\end{valuegrid}
\Zweck{valuegrid ohne Senkrechte. Vorbelegung: both.}

\begin{valuegrid}[C-22 · grid=none][grid=none]{2}
  Mustertext & zum & Formvergleich \\
  Zweite Zeile & ganz & ohne Gitter
\end{valuegrid}
\Zweck{Gar keine Linien — wo die Spalten selbst sprechen.}

\begin{defbox}[C-23 · ZSFlist ohne Titel][tone=neutral, weight=quiet]
  Ohne Titel entsteht keine Box, also gibt es keinen Träger für die
  Box-Regler. \ZSFkeyword*{ordered} und \ZSFkeyword*{tone} wirken trotzdem,
  jeder andere Regler bricht dort mit einer Meldung ab.
\end{defbox}
\Zweck{Bekannte Sonderstelle im Reglersystem, kein eigener Baustein.}
